\section{Introduction}

A 150 mm or 200 mm Dobsonian telescope shows hundreds of galaxies and clusters, yet many first-time owners see only the Moon and the planets. The mount is two bearings and a tube. It has no setting circles, the finder shows a small patch of sky, and most deep-sky objects are too faint to see in it. Observers traditionally reach them by star hopping from a bright star with a paper chart, a skill that takes many nights to learn.

Three kinds of tools address this. Digital setting circles fit encoders to both axes and report the tube's altitude and azimuth to a handset or a planetarium app such as SkySafari~\cite{skysafari}. Camera-based finders photograph the sky and identify the star pattern: Celestron's StarSense Explorer does this with the phone's own camera looking into a mirror dock~\cite{starsense}, and the open-source PiFinder uses a Raspberry Pi camera~\cite{pifinder}. The third group uses a phone strapped to the tube as the sensor. AstroHopper~\cite{astrohopper} and SkEye~\cite{skeye} read the phone's gyroscope, gravity sensor and compass, and ask the observer to point at one known star so that the app can remove the compass offset. This last approach needs no extra hardware, but the compass is unreliable next to a steel tube and the gyroscope drifts within minutes, so the observer has to realign often.

AstroFixxer started as a web port of AstroHopper and is now a native Android application. It keeps the free, sensor-based design and adds the pieces that the review in Section~\ref{sec:related} finds missing from the free tools: a large offline catalogue, an explicit and correctable alignment step, and camera plate solving that runs on the phone without a network connection or a dock.

The contributions described in this paper are:
\begin{itemize}
\item a review of free and commercial push-to tools and plate solvers, compared by the hardware they need, how they determine pointing, and whether they work offline (Section~\ref{sec:related}, Table~\ref{tab:compare});
\item an offline data pipeline that converts Stellarium's catalogues into compact Android assets, including 579,984 Gaia- and Hipparcos-based stars for plate solving (Section~\ref{sec:data});
\item a corrected coordinate model whose worst error against astropy falls from 801 to 28 arcseconds over 66 reference cases (Section~\ref{sec:core});
\item a pure-Kotlin plate solver with a statistical acceptance test, evaluated on synthetic wide-field and eyepiece images (Section~\ref{sec:solver}).
\end{itemize}

Section~\ref{sec:related} reviews related software and algorithms. Section~\ref{sec:system} describes what AstroFixxer implements, Section~\ref{sec:design} shows its design in UML diagrams, and Section~\ref{sec:review} compares it with the reviewed tools and lists what has not yet been tested. Section~\ref{sec:conclusion} concludes.
